\chapter{Infraestrutura de Montanha, Validação e Riscos}

A estação requer energia resiliente, UPS, baterias, gerador, enlaces ópticos, gestão térmica, proteção contra descargas, controle de EMC, segurança perimetral e logística. A seleção do sítio deve considerar relevo, horizonte, RFI local, clima, acesso, energia e disponibilidade de fibra. A qualificação ambiental deve abranger vento, vibração, chuva, umidade, salinidade, UV e gradientes térmicos.

O modelo de ameaça presume que todos os sítios serão eventualmente localizados, fotografados e cartografados por possíveis rivais com sensoriamento remoto óptico ou radar de abertura sintética. Sistemas SAR formam imagens de estruturas fixas a partir dos ecos de micro-ondas e operam independentemente de iluminação solar e, em grande medida, de nuvens \cite{esa_sar_intro}. Portanto, segurança e continuidade operacional não podem depender do segredo das coordenadas, da aparência visual do domo ou da expectativa de que seu eco jamais seja observado. A arquitetura deve assumir cadastro adversário prévio e apoiar-se em ausência de emissão própria, distribuição geográfica, redundância entre nós, manutenção modular, proteção física/cibernética e recuperação após perda de estação. A baixa assinatura de emissão continua útil porque dificulta inferir atividade, modo de operação e dados coletados, mesmo quando a localização física já é conhecida.

A infraestrutura temporal inclui pelo menos duas antenas GNSS externas, com coordenadas e centros de fase levantados. Elas são montadas acima ou próximas ao radome, com visão desobstruída do céu, separação física das Yagis VHF/UHF e cabos filtrados e blindados até a sala temporal. Receptores GNSS multiconstelação redundantes permitem verificação cruzada, detecção de falhas e disciplina do relógio atômico.

A interface civil proposta é uma base de concreto armado em forma de paralelepípedo, com aproximadamente $4\,\mathrm{m}\times4\,\mathrm{m}$ em planta (\texttt{CIV-001}) e $3\,\mathrm{m}$ de pé-direito estrutural (\texttt{CIV-002}). Sua laje superior deve interagir com o corte inferior proposto do radome (\texttt{GEO-007}); a geometria de apoio permanece bloqueada até C2 verificar a transição necessária ou o tamanho da base. Há uma abertura de serviço na parede frontal para permitir a entrada de pessoal e equipamentos sem remover os painéis RF.

A solução térmica é setorizada. O calor produzido por LNAs, ADCs, ASICs, FPGAs, transceptores ópticos e conversores DC/DC é conduzido para placas frias locais e rejeitado por caminhos de ar ou líquido monitorados. O núcleo central de sincronismo e fusão possui zona térmica independente. Sensores de condensação, temperatura de entrada e saída e telemetria de potência permitem controle em malha fechada. A deriva térmica de fase é medida, associada ao timestamp e incluída no modelo de calibração.

A solução estrutural utiliza uma moldura interna portante com junções vedadas, placas nodais e reforços. Cada eixo haste/lança da Yagi externa continua do vértice da pirâmide interna através da face triangular, enquanto os painéis dielétricos fornecem proteção ambiental e transmissão RF controlada. A aceitação inclui resistência estática e estabilidade dinâmica: vento, vibração, resposta modal, expansão térmica, chuva, manutenção e ciclos térmicos repetidos.

\section{Hipótese de custo--desempenho e manutenibilidade}
A modularidade por face permite ensaiar, calibrar, remover e substituir uma cadeia receptora sem desmontar uma cobertura eletromagneticamente transparente de banda ultra-ampla. A ausência de transmissor próprio também remove do nó o amplificador de potência, duplexação de alta potência e parte dos requisitos de segurança e refrigeração associados a radar ativo. Em contrapartida, antenas expostas aumentam área ao vento, risco de corrosão, gelo, descarga, dano mecânico, vedação de passagens e frequência potencial de inspeção. A arquitetura só será economicamente preferível se o custo evitado no casco e na cadeia transmissora superar esses custos adicionais e o custo da multiplicação de faces e canais.

A comparação controlada deve registrar lista de materiais, ferramental, horas de fabricação e integração, rendimento de produção, custo de instrumentação/calibração, consumo de energia, sobressalentes, tempo médio de reparo, indisponibilidade e manutenção em sítio ao longo de uma vida de referência comum. O desempenho deve ser normalizado por cobertura espectral e angular, ganho realizado, sensibilidade, precisão/covariância de localização, disponibilidade e resiliência. Sem esse modelo e cotações rastreáveis, a literatura sustenta maior complexidade de materiais e fabricação em radomes UWB multicamada, mas não prova que o RADOME tenha menor custo total.

O caminho de desenvolvimento é P0 simulação de sistema, P1 demonstrador de três nós VHF/UHF, P2 tile L/S/C, P3 tile X/Ku/Ka e P4 campanha externa integrada. Os ensaios incluem S-parâmetros, padrões 3D OTA, perda e fase do radome, coerência, sincronismo, AOA calibrado, Pd/Pfa de radar passivo e resiliência após perda de enlace ou energia.

\begin{figure}[H]
\centering
\includegraphics[width=0.95\textwidth]{pt-BR/fig10_roteiro.png}
\caption{Roteiro de desenvolvimento e validação: as cinco barras coloridas escalonadas representam gates sucessivos e parcialmente sobrepostos, desde requisitos e demonstrador VHF/UHF até tiles de faixas superiores e campanha ambiental integrada.}
\label{fig:fig10_roteiro}
\end{figure}
